% ============================================================
%  Sección 6 — Avance del proyecto, semanas 3 a 6
%  Redacción académica en lenguaje accesible.
%  Incorpora la retroalimentación del asesor sobre el Entregable 1.
% ============================================================

\section{Avance del proyecto}
\label{sec:avance}

Entre la tercera y la sexta semana se completaron los pasos 2, 3 y 4 del plan de trabajo (Tabla~\ref{tab:pasos}) y se revisaron las cifras del paso 1. El avance ya incorpora las observaciones del asesor resumidas en la Sección~\ref{sec:retro}, y los primeros apartados explican cómo se aplicó cada una.

Los resultados se obtuvieron con un programa en Python que lee el archivo exportado del sistema de compras sin modificarlo, por lo que pueden repetirse con cualquier exportación del mismo formato. Las capturas de los archivos de salida y del código están en el Anexo~\ref{sec:anexo}. Como se trabajó con pocas semanas y con etiquetas aún no confirmadas, los resultados son preliminares.

\subsection{Retroalimentación del asesor y cómo se incorpora}
\label{sec:av-retro}

El asesor señaló seis aspectos, resumidos en la Tabla~\ref{tab:av-retro}; los apartados siguientes los desarrollan.

\iffinal\else
\begin{table}[htbp]
\centering
\caption{Observaciones del asesor y su incorporación.}
\label{tab:av-retro}
\small
\begin{tabularx}{\textwidth}{@{}L{5.2cm}Y@{}}
\toprule
\textbf{Observación} & \textbf{Cómo se incorpora} \\
\midrule
Cada variable debe conocerse en el momento de emitir la alerta, antes del despacho & Se clasifican las variables según cuándo están disponibles y se excluyen del modelo las que usan información posterior (Sección~\ref{sec:av-tiempo}) \\
Distinguir una desviación respecto del historial de un error confirmado & Se adopta una terminología de tres niveles, que son desviación, caso sospechoso y error confirmado (Sección~\ref{sec:av-terminos}) \\
Distinguir el valor potencial en riesgo de una pérdida efectivamente ocurrida & La cifra de S/~3\,200 mensuales se presenta como valor potencial en riesgo, no como pérdida (Sección~\ref{sec:av-terminos}) \\
Validar las etiquetas con evidencia operativa, pues una modificación puede ser legítima & Las líneas rectificadas dejan de considerarse errores automáticamente y cada etiqueta exige una fuente documentada (Sección~\ref{sec:av-marcado}) \\
Explicar el paso del análisis semanal a alertas por pedido & Se describe el procedimiento para evaluar cada línea en el momento en que se registra (Sección~\ref{sec:av-pedido}) \\
Reservar un periodo independiente de evaluación y moderar las conclusiones & Se separan tres periodos y se reserva un extracto nuevo que no se usará hasta la evaluación final (Sección~\ref{sec:av-periodos}) \\
\bottomrule
\end{tabularx}
\end{table}
\fi

\subsubsection{Qué información existe en el momento de alertar}
\label{sec:av-tiempo}

En el archivo, cada línea se registra en promedio cinco días antes de su fecha de entrega, y el 80\,\% se registra con menos de una semana de anticipación. Esto significa que, cuando se registra una línea, todavía no se conocen las demás líneas que la misma tienda cargará esa semana, ni lo que pedirá la semana siguiente. La Tabla~\ref{tab:av-disponibilidad} explica qué significa cada variable y si está disponible en ese momento.

\begin{table}[htbp]
\centering
\caption{Variables del método, su significado y su disponibilidad al emitir la alerta.}
\label{tab:av-disponibilidad}
\footnotesize
\begin{tabularx}{\textwidth}{@{}L{2.7cm}YL{1.6cm}L{3.1cm}@{}}
\toprule
\textbf{Variable} & \textbf{Qué significa (ejemplo)} & \textbf{¿Disponible al alertar?} & \textbf{Tratamiento} \\
\midrule
Referencia de la tienda & Lo que esa tienda suele pedir de ese producto en una semana, calculado como la mediana de sus últimas cuatro semanas. \emph{Ej.: en las cuatro semanas previas al 17 de agosto, ATE pidió 216, 126, 108 y 108 panes grandes; su referencia (la mediana) es 117.} & Sí & Se calcula solo con semanas anteriores \\
Variación frente a la semana anterior & Cuánto cambió el pedido respecto de la semana pasada. \emph{Ej.: de 120 a 240 unidades es una variación de +100\,\%.} & Sí & Sin cambios \\
Línea duplicada & Indica si la tienda tiene dos líneas iguales, con el mismo producto, misma fecha de entrega y misma cantidad. Suele ser el mismo pedido cargado dos veces. \emph{Ej.: dos líneas de 180 panes junior para el 05/09.} & Sí, parcialmente & Solo se compara con las líneas ya registradas \\
Peso del producto en el pedido & Qué parte del pedido total de la tienda corresponde a ese producto, comparada con lo usual. \emph{Ej.: el pan grande suele ser el 20\,\% del pedido y esta semana es el 50\,\%.} & Parcialmente & Con lo registrado hasta ese momento \\
Presentaciones que suben juntas & Cuántas de las tres presentaciones (grande, mediano y junior) subieron más de 30\,\% en la misma tienda y semana. Si suben las tres, suele ser demanda real; si sube una sola, es más sospechoso. & Parcialmente & Con lo registrado hasta ese momento \\
Movimiento común de la cadena & Cuánto subieron o bajaron en promedio todas las tiendas esa semana. Permite distinguir un feriado, en que suben todas, de un error, en que sube una sola. \emph{Ej.: en Fiestas Patrias toda la cadena subió alrededor de 11\,\%.} & Parcialmente & Con las tiendas que ya registraron su pedido \\
Lo que la tienda pide la semana siguiente & Si el pedido alto se repite o vuelve a lo normal la semana después. Ayudó en el análisis, pero ese dato todavía no existe cuando hay que alertar. & \textbf{No} & Solo para análisis retrospectivo; fuera del modelo \\
\bottomrule
\end{tabularx}
\end{table}

Esto afecta a dos elementos del primer entregable. La última fila de la Tabla~\ref{tab:firma} usa el comportamiento de la semana siguiente; sirvió para diferenciar una campaña de un error, pero no puede usarse para alertar. Lo mismo ocurre con el parámetro de reversión de la Tabla~\ref{tab:hiper}, que queda solo para revisar casos pasados. Además, las reglas semanales de esta sección usan semanas completas, por lo que en operación se reemplazan por la versión por pedido.

\subsubsection{Desviación, caso sospechoso y error confirmado}
\label{sec:av-terminos}

En el primer entregable se usaron como sinónimos «desviación», «anomalía» y «error». A partir de este avance se usan tres términos distintos.

\begin{itemize}
\item \emph{Desviación}. Es un pedido que se aleja del patrón histórico de la tienda. No implica que esté mal; puede deberse a una promoción, a un evento local o a un cambio de necesidad.
\item \emph{Caso sospechoso}. Es una desviación que se debe revisar porque presenta al menos una de las siguientes características.
\begin{itemize}
\item la tienda subió más de 80\,\% sobre su referencia mientras el resto de la cadena no subió (no es feriado ni campaña);
\item subió una o dos presentaciones, pero no las tres a la vez, lo que no se parece a un aumento real de demanda;
\item tiene una línea duplicada, con el mismo producto, fecha de entrega y cantidad;
\item una línea es varias veces más grande que las líneas habituales de esa tienda y producto;
\item el pedido cayó más de 60\,\% respecto de su referencia, lo que puede indicar un pedido incompleto.
\end{itemize}
\item \emph{Error confirmado}. Es un caso sospechoso respaldado por evidencia operativa, por ejemplo un correo en que el área o la tienda reconoce el error.
\end{itemize}

Con esta distinción, los 17 casos y los S/~3\,200 mensuales de la Tabla~\ref{tab:exceso} corresponden a desviaciones. Esa cifra es un valor potencial en riesgo, es decir, el valor del producto despachado por encima de lo esperado si todas las desviaciones fueran errores. La pérdida real será menor y se podrá estimar cuando se sepa cuántas desviaciones se confirman y qué parte del exceso no se recuperó por otra vía.

\subsubsection{Del análisis semanal a la alerta por pedido}
\label{sec:av-pedido}

El análisis exploratorio agrupó los pedidos por tienda, producto y semana, pero la herramienta final debe evaluar cada línea cuando se registra. Para ello se combinan tres comparaciones.

\begin{enumerate}
\item La referencia de la tienda sigue siendo semanal y se calcula con las semanas anteriores, que ya se conocen.
\item Al registrarse una línea, se suma con las líneas de la misma tienda y producto ya cargadas para esa semana, y el acumulado se compara con lo que la tienda suele llevar pedido hasta ese día.
\item La línea se compara también con el tamaño habitual de las líneas de esa tienda y producto, para detectar errores que en el total semanal no se notan (Sección~\ref{sec:av-marcado}).
\end{enumerate}

Así, la alerta corresponde a una línea concreta y puede emitirse el mismo día del registro, varios días antes de la entrega.

\subsubsection{Periodos de ajuste, validación y prueba}
\label{sec:av-periodos}

Hasta ahora el extracto se dividía en dos partes, hasta el 10 de agosto para ajustar y las tres semanas siguientes para evaluar. Sin embargo, esas semanas también se revisaron al definir los casos sospechosos, así que no eran del todo independientes. Se adopta entonces una división en tres periodos, que deja las dos últimas semanas completas sin usar hasta el final.

\begin{itemize}
\item Ajuste (13 de julio al 3 de agosto), para construir el método y fijar sus umbrales.
\item Validación (10 y 17 de agosto), para comparar alternativas y elegir parámetros.
\item Prueba (24 y 31 de agosto), que se mide una sola vez, con el método ya definido.
\end{itemize}

La semana del 7 de septiembre está incompleta y no se usa. Además, se pedirá al área de planeamiento un extracto posterior al 11 de septiembre para confirmar los resultados con pedidos nuevos. La precisión y la exhaustividad se medirán en el periodo de prueba cuando haya etiquetas confirmadas.

\subsection{Resumen del estado del proyecto}
\label{sec:av-resumen}

El proyecto se encuentra en la semana 6 de 12, y los productos comprometidos hasta esta semana están terminados. Lo que todavía no se puede afirmar es si las alertas alcanzarán la precisión buscada, porque falta confirmar los casos de error. La Tabla~\ref{tab:av-estado} resume la situación.

\begin{table}[htbp]
\centering
\caption{Estado del proyecto al cierre de la semana 6 de 12. La columna «Semana» indica cuándo debe cumplirse cada meta. La precisión y la exhaustividad se medirán en el periodo de prueba cuando haya etiquetas confirmadas.}
\label{tab:av-estado}
\small
\begin{tabularx}{\textwidth}{@{}YL{1.9cm}cL{2.3cm}L{2.2cm}@{}}
\toprule
\textbf{Aspecto evaluado} & \textbf{Meta} & \textbf{Semana} & \textbf{Situación en la semana 6} & \textbf{Estado} \\
\midrule
Productos de los pasos 1 a 4 & 4 de 4 & 6 & 4 de 4 & \textcolor{green!50!black}{\textbf{Cumplido}} \\
Errores confirmados con evidencia operativa & 20 o más & 7 & 1 (caso ATE) & \textcolor{orange!85!black}{\textbf{En camino}} \\
Alertas por semana (acumulado de la semana) & Entre 2 y 5 & 8 & 4,0 & \textcolor{green!50!black}{\textbf{En rango}} \\
Alertas que coinciden con un caso sospechoso & 70\,\% o más & 12 & Por medir & \textcolor{orange!85!black}{\textbf{Pendiente}} \\
Casos sospechosos que el método detecta & 80\,\% o más & 12 & Por medir & \textcolor{orange!85!black}{\textbf{Pendiente}} \\
\bottomrule
\end{tabularx}
\end{table}

Por ahora solo hay un error confirmado por escrito (ATE); los demás son casos sospechosos. Por eso, en esta etapa los métodos se prueban con casos conocidos (Sección~\ref{sec:av-casos}) y la medición de precisión y exhaustividad queda para cuando haya más errores confirmados.

\subsection{Revisión de la limpieza de datos}
\label{sec:av-limpieza}

El programa de limpieza reprodujo casi todas las cifras de la Sección~\ref{sec:anuladas}, entre ellas 10\,374 líneas vigentes, 385 anuladas con reemplazo y 119 líneas duplicadas. Cambiaron dos valores.

Las líneas anuladas que se reemplazaron por otra cantidad pasan de 27 a 23. Cuando una tienda tiene varias líneas del mismo producto el mismo día, no es evidente cuál reemplazó a cuál; el programa toma la primera según el número de orden y posición.

Las observaciones evaluables bajan de 1\,554 a 1\,371 porque se excluyó la última semana. El archivo termina un viernes y esa semana no tiene los pedidos del sábado, por lo que aparentaba una caída general.

\subsection{Paso 2: la tabla de señales}
\label{sec:av-senales}

Para cada tienda, producto y semana se calcularon las ocho señales de la Tabla~\ref{tab:variables}. Estas indican cuánto se aleja el pedido de lo que la tienda suele pedir, cuánto cambió frente a la semana anterior, si hay una línea repetida, si subieron las tres presentaciones y si toda la cadena se movió en la misma dirección. Esta versión usa semanas completas; la versión por pedido se describe en la Sección~\ref{sec:av-por-pedido}.

\begin{figure}[htbp]
\centering
\includegraphics[width=\textwidth]{xl_02_senales.png}
\caption{Captura de \texttt{salidas/02\_senales.xlsx} abierto en Excel, con una fila por tienda, producto y semana con sus señales. El procedimiento que la genera está en el Algoritmo~\ref{alg:senales}.}
\label{fig:av-cap-senales}
\end{figure}

La tabla tiene 1\,371 filas evaluables. La semana del 27 de julio vuelve a mostrar un alza de toda la cadena, que coincide con Fiestas Patrias. La semana del 13 de julio muestra el mismo comportamiento sin un motivo conocido, y se consultará al área de planeamiento. El pedido de ATE del 17 de agosto vuelve a aparecer marcado.

\subsection{Paso 3: casos sospechosos y su validación}
\label{sec:av-marcado}

El archivo no indica qué pedidos estuvieron mal, así que la lista de errores se construye en dos pasos. Primero, el programa selecciona pedidos sospechosos y propone una etiqueta. Después, cada propuesta se confirma o descarta con evidencia operativa.

Se seleccionaron 78 pedidos por tres motivos, que son un alza solo en esa tienda, una línea duplicada o una caída brusca. Se propone como posible error cuando la tienda subió sola, sin que subieran las demás tiendas ni sus otras presentaciones, o cuando hay una línea duplicada con un alza importante (Tabla~\ref{tab:av-candidatos}).

\begin{table}[htbp]
\centering
\caption{Pedidos sospechosos según el motivo y la etiqueta propuesta.}
\label{tab:av-candidatos}
\small
\begin{tabular}{@{}lrr@{}}
\toprule
\textbf{Motivo} & \textbf{Posible error} & \textbf{Probablemente normal} \\
\midrule
Alza que solo afectó a la tienda & 8 & 8 \\
Línea duplicada & 5 & 38 \\
Caída brusca & 0 & 19 \\
\midrule
\textbf{Total} & \textbf{13} & \textbf{65} \\
\bottomrule
\end{tabular}
\end{table}

\begin{figure}[htbp]
\centering
\includegraphics[width=\textwidth]{xl_03_marcado.png}
\caption{Hoja de revisión \texttt{salidas/03\_marcado.xlsx}. Cada fila es un caso sospechoso con su motivo, lo habitual de la tienda y los identificadores de sus líneas. En rojo, los casos propuestos como posible error; en la columna \texttt{etiqueta} la autora escribe \emph{error}, \emph{normal} o \emph{duda}, y en \texttt{evidencia}, la fuente. El caso C11 (ATE) ya figura como error confirmado.}
\label{fig:av-cap-marcado}
\end{figure}

La autora valida cada caso y lo marca como error solo si cuenta con alguna de estas evidencias, como un correo en que el área o la tienda reconozca el error, una corrección del pedido con su motivo o la confirmación de la analista que atendió el caso. Los casos sin evidencia quedan como «duda» y no se usan para entrenar ni para evaluar. La forma de registrar estas etiquetas se describe en la Sección~\ref{sec:etiqueta}.

En el primer entregable se propuso usar como ejemplos de error las líneas que compras anuló y volvió a cargar con otra cantidad. Esta idea se descartó. Como indicó el asesor, cambiar una cantidad puede ser una decisión legítima de la tienda. Además, en 19 de las 23 líneas la cantidad corregida es mayor que la original (por ejemplo, de 18 a 72 unidades), lo que puede deberse tanto a un error de digitación como a que la tienda decidió pedir más. Estas líneas se tratarán como casos sospechosos que requieren evidencia.

También se observó que ninguna de estas correcciones sobresale en el total semanal, porque afectan a una sola línea. Por eso se agregó la comparación por línea descrita en la Sección~\ref{sec:av-pedido}.

\subsection{Paso 4: reglas sencillas como referencia}
\label{sec:av-reglas}

Antes de entrenar el árbol de decisión se programaron tres reglas sencillas, que sirven de referencia para comparar el árbol.

\begin{enumerate}[label=\textbf{R\arabic*.}]
\item Semana anterior. Alerta si el pedido duplica el de la semana previa. Es lo más parecido a lo que hoy hace la analista.
\item Habitual de la tienda. Alerta si el pedido se aleja de lo que la tienda suele pedir más de lo que normalmente varía.
\item Inusual sin ejemplos. Es un método llamado Isolation Forest \cite{liu2008}, que no usa las etiquetas. Construye muchos árboles que parten los datos al azar y mide cuántas divisiones hacen falta para dejar aislado cada pedido. Un pedido muy distinto del resto queda aislado con pocas divisiones y se marca como inusual. Su limitación es que detecta lo \emph{raro}, no lo \emph{erróneo}, ya que una tienda grande o una semana de feriado también son raras sin ser errores.
\end{enumerate}

Las dos primeras reglas usan umbrales fijos y la tercera se ajustó con el periodo de ajuste. Como todavía hay un solo error confirmado, en esta etapa no se reporta su precisión, porque con tan pocos casos la cifra no sería concluyente. Se analizan, en cambio, dos aspectos que no dependen de las etiquetas, es decir, cuántas alertas emite cada regla por semana (Figura~\ref{fig:av-alertas}) y cómo se comporta frente a dos casos conocidos (Sección~\ref{sec:av-casos}).

\begin{figure}[htbp]
\centering
\includegraphics[width=0.85\textwidth]{fig_06_alertas_semana.png}
\caption{Número de alertas por semana de cada regla. La franja sombreada marca el volumen que el analista puede revisar; las líneas verticales marcan el inicio de la validación y de la prueba.}
\label{fig:av-alertas}
\end{figure}

La regla de la semana anterior y la del método sin ejemplos se mantienen casi siempre dentro de la franja revisable. La regla de lo habitual de la tienda emite más alertas, y en la semana de Fiestas Patrias llega a 34, porque no diferencia una tienda que sube sola de toda la cadena subiendo por el feriado. El árbol de decisión podría corregir esto, ya que dispone de la señal de movimiento común.

\subsection{Primera versión por pedido}
\label{sec:av-por-pedido}

Se programó una primera versión que evalúa cada línea en el momento de su registro, con las tres comparaciones de la Sección~\ref{sec:av-pedido}. Cada línea se compara solo con las semanas anteriores y con las líneas de la misma semana ya registradas. Los umbrales se fijaron con el periodo de ajuste, buscando unas cuatro alertas por semana. Con ellos, la comparación del acumulado de la semana emitió 4,0 alertas por semana en ese periodo, la del tamaño de la línea 3,3 y la combinación de ambas con la línea duplicada 13,5.

\begin{figure}[htbp]
\centering
\includegraphics[width=\textwidth]{xl_05_alertas.png}
\caption{Captura de la hoja \texttt{alertas} de \texttt{salidas/05\_alertas\_pedido.xlsx}. Cada fila es una alerta emitida en la primera línea que la dispara, con su fecha de registro, su fecha de entrega y los días de anticipación.}
\label{fig:xl-alertas}
\end{figure}

Como cada alerta se emite cuando se registra la línea, también se conoce con cuánta anticipación llega. En el caso de ATE, la línea de 396 unidades se registró el 19 de agosto para entregarse el 22, y la alerta se habría emitido tres días antes del despacho.

\subsection{Primer entrenamiento del árbol y análisis de parámetros}
\label{sec:av-arbol}

Con las señales por línea y las etiquetas propuestas se entrenó una primera versión del árbol, conectada a la ventana descrita en la Sección~\ref{sec:arquitectura}. Para analizar el efecto de sus parámetros se entrenaron 24 combinaciones con el periodo de ajuste y se comparó cuántas alertas emite cada una y con qué precisión. La Tabla~\ref{tab:av-parametros} muestra las más representativas.

\begin{table}[htbp]
\centering
\caption{Análisis de parámetros del árbol entrenado con el periodo de ajuste (33 líneas de error entre 3\,902). Las alertas se cuentan por tienda, producto y semana.}
\label{tab:av-parametros}
\small
\begin{tabular}{@{}lcccc@{}}
\toprule
\textbf{Peso de la clase error} & \textbf{Profundidad} & \textbf{Hojas} & \textbf{Alertas por semana} & \textbf{Precisión} \\
\midrule
Balanceado & 2 & 4 & 42,5 & 2\,\% \\
Balanceado & 5 & 9 & 37,8 & 3\,\% \\
1 (sin ajuste) & 2 a 5 & 4 a 9 & 1,0 & 75\,\% \\
\textbf{5} & \textbf{4} & \textbf{7} & \textbf{1,8} & \textbf{57\,\%} \\
5 (mín. 20 líneas por hoja) & 4 & 5 & 7,5 & 13\,\% \\
\bottomrule
\end{tabular}
\end{table}

La tabla muestra el efecto de cada parámetro. Con el peso balanceado, que iguala la importancia total de errores y normales, el árbol emite unas 40 alertas por semana, demasiadas para revisarlas. Sin ajuste de peso ocurre lo contrario, ya que casi no alerta. Con un peso de 5 y profundidad 4 la carga baja a menos de 2 alertas por semana con 57\,\% de precisión en ajuste. Exigir al menos 20 líneas por hoja vuelve las reglas más generales, pero sube la carga y baja la precisión.

Como en el periodo de ajuste hay solo 4 casos de error, estas cifras describen cómo responde el árbol a cada parámetro, pero todavía no sirven para elegir el mejor por su capacidad de detectar errores nuevos. Por eso los parámetros se fijaron con dos criterios, una carga de alertas revisable y reglas fáciles de leer (profundidad 4, mínimo 5 líneas por hoja y peso 5 para la clase error). Con esos parámetros, el árbol se entrenó con los periodos de ajuste y validación; el periodo de prueba quedó reservado.

\begin{figure}[htbp]
\centering
\includegraphics[width=0.75\textwidth]{app_entrenar.png}
\caption{Ventana de entrenamiento del prototipo, con los datos usados, cantidad de errores etiquetados y reglas aprendidas por el árbol.}
\label{fig:app-entrenar}
\end{figure}

\begin{figure}[htbp]
\centering
\includegraphics[width=\textwidth]{app_revisar.png}
\caption{Ventana de revisión del prototipo, con las líneas marcadas por el árbol en las últimas semanas del extracto, con su nivel, la probabilidad de la hoja y la regla en palabras.}
\label{fig:app-revisar}
\end{figure}

\subsection{Prueba con casos conocidos}
\label{sec:av-casos}

Mientras no haya suficientes errores confirmados, los métodos se prueban con dos casos cuyo resultado se conoce de antemano y que el primer entregable ya comparó en la Tabla~\ref{tab:firma}.

\begin{itemize}
\item \textbf{ATE, semana del 17 de agosto}. Es un error confirmado por correo. Un buen método \emph{debe} alertar.
\item \textbf{Fiestas Patrias, semana del 27 de julio}. Es un alza de toda la cadena, que no es error. Un buen método \emph{no debe} llenarse de alertas; esa semana hay 169 combinaciones de tienda y producto.
\end{itemize}

Para que la prueba sea justa, el árbol se entrenó sin el caso que se evalúa, es decir, sin la semana de ATE para ver si la detecta, y sin la semana de Fiestas Patrias para contar sus alertas. La Tabla~\ref{tab:av-casos} resume el resultado.

\begin{table}[htbp]
\centering
\caption{Prueba con casos conocidos. En Fiestas Patrias se cuentan las alertas emitidas, que en ese caso son falsas alarmas.}
\label{tab:av-casos}
\small
\begin{tabularx}{\textwidth}{@{}Ycc@{}}
\toprule
\textbf{Método} & \textbf{ATE, 17 de agosto (debe alertar)} & \textbf{Fiestas Patrias (alertas)} \\
\midrule
R1. Semana anterior & Sí & 3 \\
R2. Habitual de la tienda & Sí & 34 \\
R3. Inusual sin ejemplos & Sí & 3 \\
\midrule
Por pedido, acumulado de la semana & \textbf{Sí, 3 días antes} & \textbf{0} \\
Por pedido, tamaño de la línea & Sí, 3 días antes & 9 \\
Por pedido, cualquiera o línea duplicada & Sí, 3 días antes & 16 \\
\midrule
Árbol de decisión (sin haber visto el caso) & No & 3 \\
\bottomrule
\end{tabularx}
\end{table}

La comparación del acumulado de la semana es la única que cumple ambas condiciones, ya que detecta el error de ATE tres días antes de la entrega y no emite ninguna alerta en Fiestas Patrias, porque descuenta el movimiento común de la cadena. La regla de lo habitual de la tienda detecta ATE, pero se llena de falsas alarmas en el feriado. El árbol, en cambio, distingue bien el feriado, pero no reconoce el caso de ATE cuando no lo vio durante el entrenamiento. Esto confirma que, con los pocos errores etiquetados que hay hoy, el árbol todavía no aprende un patrón que pueda aplicar a casos nuevos.

\subsection{Qué hace falta para medir el desempeño}
\label{sec:av-mejoras}

La precisión y la exhaustividad, que son las metas del proyecto, se medirán una sola vez en el periodo de prueba (24 y 31 de agosto), que quedó reservado. Hacerlo hoy no sería concluyente por tres razones. Hay un solo error confirmado. En esas dos semanas hay apenas 5 casos propuestos, de modo que cada acierto movería el resultado 20 puntos. Además, las etiquetas propuestas fueron definidas por el mismo programa. Para que esa medición sea válida hace falta lo siguiente, en este orden.

\begin{enumerate}
\item \textbf{Ampliar el historial} con el extracto posterior al 11 de septiembre y, si es posible, con meses anteriores a julio, para tener más casos de error de los cuales aprender y con los cuales evaluar.
\item \textbf{Medir la versión por pedido contra etiquetas por línea}, que es la forma en que esa versión emite sus alertas.
\item \textbf{Combinar la regla de lo habitual con el filtro de campaña y de presentaciones}, y dar al árbol esa medida como una variable más, para que pueda aprovechar lo que hoy hace bien cada método.
\end{enumerate}

Ninguna de estas mejoras se ajustará mirando el periodo de prueba, para que su resultado siga siendo una medición independiente.

\iffinal\else
\subsection{Riesgos observados durante la ejecución}
\label{sec:av-riesgos}

La ejecución de estas semanas y la revisión del asesor permitieron actualizar los riesgos de la Tabla~\ref{tab:riesgos} (Tabla~\ref{tab:av-riesgos}). La responsable de todas las respuestas es la autora del proyecto.

\begin{table}[htbp]
\centering
\caption{Riesgos actualizados al cierre de la semana 6.}
\label{tab:av-riesgos}
\small
\begin{tabularx}{\textwidth}{@{}L{3.6cm}L{3.2cm}L{1.6cm}Y@{}}
\toprule
\textbf{Riesgo} & \textbf{Señal observada} & \textbf{Estado} & \textbf{Respuesta} \\
\midrule
Si no se reúnen suficientes errores confirmados, el modelo no tendrá de qué aprender & Un solo error confirmado por escrito & Ocurriendo & Revisar los 78 casos contra los correos en la semana 7; si no se alcanzan 20, reportar el método como apoyo a la revisión y no como clasificador \\
Si el modelo usa información que no existe al alertar, su desempeño real será menor que el medido & Variables calculadas sobre semanas completas & Identificado & Recalcular las señales con lo registrado hasta cada momento (Sección~\ref{sec:av-pedido}) \\
Si las etiquetas se construyen con la misma idea que una regla, esa regla parecerá mejor de lo que es & La segunda regla rinde mucho más en validación que en ajuste & Vigente & Basar la etiqueta final en evidencia operativa, no en la regla \\
Si no se obtiene un extracto posterior, no habrá periodo de prueba independiente & El extracto actual termina el 11 de septiembre & Nuevo & Solicitar el extracto de septiembre y octubre al área de planeamiento en la semana 7 \\
Si la semana del 13 de julio no fue una campaña, el método podría silenciar errores en semanas similares & Movimiento conjunto sin feriado conocido & Nuevo & Consultar por escrito al área de planeamiento \\
\bottomrule
\end{tabularx}
\end{table}

\subsection{Plan para las semanas 7 a 9}
\label{sec:av-plan}

Los hitos del plan general se mantienen, pero las tareas de las próximas semanas se ajustan para responder a la retroalimentación (Tabla~\ref{tab:av-plan}).

\begin{table}[htbp]
\centering
\caption{Tareas de las semanas 7 a 9 y sus dependencias.}
\label{tab:av-plan}
\small
\begin{tabularx}{\textwidth}{@{}cYcL{3cm}c@{}}
\toprule
& \textbf{Tarea} & \textbf{Días} & \textbf{Requiere terminar antes} & \textbf{Semana} \\
\midrule
A & Solicitar el extracto posterior al 11 de septiembre & 1 & -- & 7 \\
B & Revisar los 78 casos sospechosos contra los correos y registrar la evidencia & 3 & -- & 7 \\
C & Recalcular las señales con la información disponible en cada momento y añadir la comparación por línea & 3 & -- & 7 \\
D & Entrenar el árbol y compararlo con la mejor regla en validación & 2 & B y C & 8 \\
\multicolumn{5}{@{}l}{\quad\textit{Hito: decidir si el árbol se justifica. Decide la autora con el visto bueno del asesor.}} \\
E & Elegir a partir de qué desviación se emite una alerta y redactar los motivos & 2 & D & 8 \\
F & Construir el motor de cálculo que evalúa cada línea registrada & 4 & E & 9 \\
G & Evaluar una sola vez sobre el extracto independiente & 1 & A y F & 9 \\
\bottomrule
\end{tabularx}
\end{table}

Respecto del plan original, el hito del paso 5 se desplaza de la semana 7 a la 8, porque la revisión con evidencia y el recálculo de las señales deben completarse antes de entrenar el árbol. Este desplazamiento se absorbe dentro de la semana 8 y no afecta la entrega de la semana 12. La dependencia más frágil es la revisión de casos (B), porque depende de los correos que conserva el área de planeamiento; si se retrasa, la elección del umbral (E) se simplificará comparando solo tres valores posibles.

\fi
